% COMPOSE-FM -- NeurIPS workshop format (4 pages + refs)
% RULE 0: this is a FLOW MATCHING METHOD paper. Finance and biology are validation.
% The method section must contain no domain vocabulary.
% NOTE ON BUILDING: no LaTeX toolchain was available in the environment where this was
% written, so this file has NOT been compiled. Drop `neurips_2025.sty` (or the style for
% the target workshop) next to it and build with pdflatex; if the style file is absent the
% \IfFileExists fallback below gives a plain-article build with the right geometry so the
% content can be read and page-counted.
\documentclass{article}
\IfFileExists{neurips_2025.sty}{\usepackage[final]{neurips_2025}}{%
  \usepackage[margin=1in]{geometry}\usepackage{natbib}}
\usepackage[utf8]{inputenc}
\usepackage[T1]{fontenc}
\usepackage{hyperref}
\usepackage{url}
\usepackage{booktabs}
\usepackage{amsfonts,amsmath,amssymb,amsthm}
\usepackage{graphicx}
\usepackage{xcolor}

\newtheorem{proposition}{Proposition}
% Figure path indirection. Defined once here so the include site below has unambiguous
% single-brace syntax: an artifact marker written directly inside \includegraphics{...}
% yields {{{...}}}, which graphicx cannot parse.
%
% This is a LaTeX source file meant to be compiled with pdflatex, so the path is a plain
% relative path and the figure must sit at paper/figures/fig2_case_studies.png (it does,
% in the repo). Do not substitute a platform artifact marker here: nothing expands it at
% compile time, so it would reach graphicx verbatim and fail. The figure is also stored as
% an artifact (fig2_case_studies.png, artifact_id
% e4061ed5-fad8-4faf-9652-a2dba830d479, version_id
% 511f4f32-ae16-4029-b367-617206186d6b) if you need to retrieve the image itself.
\newcommand{\figcases}{figures/fig2_case_studies.png}

\newcommand{\XP}{X_{\mathcal{P}}}
\newcommand{\Rd}{\mathbb{R}^d}

\title{Compose the Generators, Not the Displacements:\\
Covariant Composition of Intervention Flows}

\author{%
  Anonymous Author(s)\\
  Affiliation\\
  \texttt{email}
}

\begin{document}
\maketitle

\begin{abstract}
Models that predict the effect of combined interventions in a latent space almost
universally \emph{add displacements}: each intervention is mapped to a vector, and a
combination is the sum. We show this is not a chart-independent operation. Under a
diffeomorphism of the latent space, a sum of displacements transforms with an error of
order $\|d\|^2$, so the prediction depends on an arbitrary choice of coordinates, whereas
a sum of \emph{vector fields} is exactly covariant. We measure the defect: summing fields
reproduces the reference to $5.5\times 10^{-17}$ while summing displacements errs by
$1.7\times 10^{-3}$, and the error grows with the number of interventions, the exposure,
and the anisotropy of the latent geometry. This motivates COMPOSE-FM, a flow-matching
operator that composes learned \emph{generators}: a state-dependent contraction gate, a
sparse pairwise coupling term made tensorial by a Levi-Civita connection correction, and
exposure entering as integration time, which makes the exposure semigroup exact by
construction rather than fitted. On two unrelated domains -- limit-order-flow composition
in cryptocurrency markets and a patient-derived organoid drug-combination ladder with a
held-out triple -- composing generators improves over a trained displacement baseline of
matched capacity by $2.9\times$ and $3.0\times$ respectively on intervention sets never
seen during training. We also report where the method does not help, and why.
\end{abstract}

\section{Introduction}

Predicting the effect of applying several interventions at once is a recurring problem:
which perturbations combine super-additively in a cell, how simultaneous order flow moves
a price, how two treatments interact. A dominant modelling pattern encodes each
intervention as a latent displacement and predicts a combination by adding those
displacements. This pattern underlies compositional perturbation autoencoders and related
latent perturbation models, and it is attractive because it is simple and because the
additive prediction is a sensible null.

Our starting point is that this pattern has a geometric defect that has not, to our
knowledge, been stated. A latent space learned by an autoencoder or by PCA is defined only
up to a smooth reparameterisation: nothing in the training objective distinguishes $z$ from
$\psi(z)$ for a diffeomorphism $\psi$. A modelling operation that is meant to express a
fact about interventions should therefore not depend on which of those equivalent charts
was chosen. Adding displacements does depend on it. Adding \emph{vector fields} does not.

This is not a philosophical point; it is measurable, and it grows in exactly the regimes
practitioners care about -- more simultaneous interventions, stronger exposures, more
anisotropic latent geometry (\S\ref{sec:defect}). It also suggests what to replace the
pattern with: a model whose primitive object is a generator $X_p$ rather than a
displacement $d_p$, and whose composition rule acts on fields.

\paragraph{Contributions.}
\begin{enumerate}
\item We identify and measure a covariance defect in displacement-additive composition,
      and show a sum of generators is exactly covariant while a sum of displacements is
      not (\S\ref{sec:defect}).
\item We give COMPOSE-FM, a flow-matching composition operator with a state-dependent
      contraction gate, a saturating exposure response with a \emph{learned} exponent, and
      a sparse pairwise coupling term that we make tensorial with a connection correction
      (\S\ref{sec:method}). Exposure is integration time, so the exposure semigroup holds
      to solver accuracy ($2.3\times10^{-7}$) rather than being learned.
\item We validate on two unrelated domains and report the negatives as carefully as the
      positives (\S\ref{sec:experiments}).
\end{enumerate}

\section{The covariance defect}
\label{sec:defect}

Let $\psi:\Rd\to\Rd$ be a diffeomorphism, understood as a change of latent chart. A
vector field transforms as $\tilde X(\tilde z) = D\psi(z) X(z)$, so a sum of fields
transforms as the sum -- composition of generators commutes with the change of chart
exactly. A displacement is not a vector at a point but a difference of two points, and
$\psi(z + d) - \psi(z) = D\psi(z)d + O(\|d\|^2)$: the quadratic remainder is precisely the
chart dependence.

We measure this on a controlled construction where the ground truth is known. Taking two
generators, a random polynomial diffeomorphism, and comparing the pushforward of the
composed prediction against the prediction composed natively in the new chart, we find

\begin{center}
\begin{tabular}{lc}
\toprule
composition rule & relative discrepancy \\
\midrule
sum of vector fields & $5.5\times10^{-17}$ \\
integrate the summed field & $3.5\times10^{-14}$ \\
sum of displacements & $1.7\times10^{-3}$ \\
\bottomrule
\end{tabular}
\end{center}

\noindent The first two are at machine and solver precision; the third is not, and it is
the operation that displacement-additive models perform. The defect is not a fixed
constant: it grows along every axis that makes the problem harder. Increasing the number
of simultaneous interventions from $k{=}2$ to $k{=}5$ raises it from $2.79\times10^{-2}$ to
$5.97\times10^{-2}$; increasing exposure $\tau$ from $0.5$ to $3.0$ raises it from
$1.40\times10^{-2}$ to $8.19\times10^{-2}$; increasing latent anisotropy from $1$ to $30$
raises it from $2.79\times10^{-2}$ to $3.98\times10^{-2}$. In a jointly reachable regime
($k{=}4$, $\tau{=}2$, anisotropy $10$) the displacement error reaches $7.84\times10^{-2}$,
$12.1\times$ the measurement noise floor, while generator composition stays below
$10^{-8}$ throughout.

\paragraph{Honest boundary.} At $k{=}2$, $\tau{=}1$ and generic (near-isotropic) geometry
the defect sits at $1.36\times$ the noise floor -- detectable but not dominant. The defect
matters for many interventions, strong exposures, or curved latent geometry; it does not
by itself explain a failure at $k{=}2$.

\section{COMPOSE-FM}
\label{sec:method}

\paragraph{Setup.} A state is $z\in\Rd$. Each intervention $p$ carries a learned autonomous
generator $X_p(z)$ and an exposure $\tau_p\ge 0$, the parameter of the one-parameter family
generated by $X_p$. For a set $\mathcal{P}$ the operator is
%
\begin{equation}
\XP(z) \;=\; \underbrace{\gamma(z,\mathcal{P})\sum_{p\in\mathcal{P}}\tau_p X_p(z)}_{\text{gated additive part}}
\;+\; \underbrace{\tfrac12\sum_{p<q}\beta_{pq}(z)\,S^{\nabla}_{pq}(z)}_{\text{coupling}},
\label{eq:operator}
\end{equation}
%
and the prediction is the time-one flow $z(1)=\phi^1_{\XP}(z(0))$, trained by flow
matching against the observed post-intervention distribution. Every term below is a
function of the state and the intervention set only; no domain quantity appears.

\paragraph{Contraction gate.} Interventions act through shared downstream capacity, so
their joint effect generally falls short of the sum. We take
$\gamma(z,\mathcal{P}) = g(z,\mathcal{P})\,\bigl(1+(T/\kappa(z))^{\nu}\bigr)^{-1}$ with
$T=\sum_p \tau_p$, where $g\in(0,1)$ is a permutation-invariant function of the state and
the active generators, and $\kappa>0,\nu>0$ are learned. The second factor is what makes
exposure extrapolation possible: a purely multiplicative gate still leaves the drive
growing linearly in $T$, so $\|\XP\|\sim T^{1-\nu}$ for $T\gg\kappa$ lets the operator
represent a saturating response law with a \emph{learned} exponent rather than an assumed
one. We do not assume $\nu=1/2$; we fit it.

\paragraph{Coupling, and why it needs a connection.} The natural second-order object is the
symmetric field $S_{pq}=DX_q X_p + DX_p X_q$, which is the leading term in
joint-minus-additive. But $S_{pq}$ is \emph{not} a tensor: measured pushforward-versus-native
discrepancy is $1.0\times10^{-1}$, i.e. the same defect one derivative up. Subtracting the
Levi-Civita connection of a learned metric,
$S^{\nabla}_{pq} = DX_q X_p + DX_p X_q - 2\Gamma(z)[X_p,X_q]$,
restores tensoriality to $1.2\times10^{-11}$. (The antisymmetric part, the Lie bracket, is
already tensorial at $1.7\times10^{-16}$, which is why bracket-based methods never meet
this problem -- and why the symmetric part, which carries the magnitude of the
interaction, does.) The weights $\beta_{pq}(z)$ use a hard-concrete $L_0$ gate so that
exact zeros are representable: the model can state that a pair does not interact in a given
state.

\paragraph{Exposure is integration time.} Because each $X_p$ is autonomous, integrating a
fixed field for time $s+t$ equals integrating for $s$ then $t$. We verify this on the
trained model at $2.3\times10^{-7}$ relative error, i.e. exact up to the RK2 solver. We
state the property precisely: it holds for the exposure-parameterised \emph{flow}. The
weaker claim ``scaling $\tau$ by $c$ equals scaling time by $c$'' is only approximate
($3.6\times10^{-5}$) because $\gamma$ depends on $\tau$, and we do not claim it.

\paragraph{Implementation.} The coupling requires Jacobian-vector products for every pair
at every integration substep. Batching all pairs into a single JVP call makes the cost
independent of $k$ (exact to $0.0$ against the naive loop, $\approx 2\times$ faster), and
computing the connection's third term as one reverse-mode gradient of a scalar rather than
$d$ forward JVPs reduces that term from $O(d)$ passes to $O(1)$.

\section{Experiments}
\label{sec:experiments}

We evaluate with normalised energy distance,
$\mathrm{ED}(\hat\mu,\mu)/\mathrm{ED}(\mu_0,\mu)$, where $\mu_0$ is the pre-intervention
distribution: $1.0$ means no better than predicting no change, $0$ is perfect. The
baseline is a trained displacement model of matched trunk width, depth and optimisation
budget, differing only in the composition rule. Populations are unpaired, so all models are
trained with a distributional (sliced-$W_2$ plus mean) objective; a squared-error objective
drives predictions to the conditional mean, which this metric correctly punishes more than
inaction.

\begin{figure}[t]
\centering
\includegraphics[width=\textwidth]{\figcases}
\caption{Composing generators generalises to intervention sets never seen in training, in
two unrelated domains. \textbf{(a)} Order-flow composition, $6{,}373$ held-out windows with
all four channels active. \textbf{(b)} Organoid drug combinations, held-out triple
Cetuximab $+$ SN-38 $+$ 5-FU, $26$ populations never seen jointly; here the no-change
baseline coincides with the best fitted-mean-shift oracle. \textbf{(c)} The exposure
response is sub-linear in the data: binning by executed quantity and fitting on bin means
gives exponent $0.59$ over a $1{,}547\times$ span of volume, which is why the gate is given
a saturating form with a learned exponent. Hatched bars mark the saturating variant.}
\label{fig:cases}
\end{figure}

\paragraph{Domain 1: order-flow composition.} From four days of BTCUSDT aggregated trades
($3.5$M trades) we build $34{,}545$ ten-second windows. The state is six microstructure
features including the realised return; the interventions are four order-flow channels
(side $\times$ size class); exposure is executed quantity normalised to each channel's
median. Training uses windows with at most three channels active; evaluation uses windows
with all four. Composing generators reaches $0.691$ against $2.008$ for the displacement
baseline ($2.9\times$), with each state-dependent rung contributing
(Fig.~\ref{fig:cases}a).

\paragraph{Domain 2: organoid drug combinations.} We use the Trellis mass-cytometry
organoid screen \citep{ramoszapatero2023}, in the preprocessed form released with Meta Flow
Matching \citep{atanackovic2024mfm}: $24{,}756{,}030$ cells $\times$ $44$ markers, $666$
populations with controls matched per condition. The ladder contains six single compounds,
four pairs and one triple. We hold out the triple entirely. Composing generators reaches
$0.741$ against $2.233$ ($3.0\times$), and the state-gated variants are the only models
that beat the fitted-oracle reference of $1.000$ (Fig.~\ref{fig:cases}b).

\paragraph{Replication across seeds.} Single-seed comparisons are not a result. Repeating
the decisive contrast over three seeds (Table~\ref{tab:seeds}), every composition variant
beats the displacement baseline in $3/3$ seeds in \emph{both} domains, with non-overlapping
spreads; the mean-over-seeds improvement is $3.1\times$ in finance and $3.3\times$ on the
organoid triple.

\begin{table}[t]
\centering
\caption{Three-seed replication on the held-out composition axis (normalised energy
distance, mean $\pm$ sd over seeds $0,1,2$; lower is better). ``Beats'' counts seeds in
which the model scores below the displacement baseline trained with the same seed. Numbers
from \texttt{results/seed\_robustness.json}.}
\label{tab:seeds}
\begin{tabular}{lcccc}
\toprule
& \multicolumn{2}{c}{finance (all 4 channels)} & \multicolumn{2}{c}{organoid (held-out triple)} \\
\cmidrule(lr){2-3}\cmidrule(lr){4-5}
& score & beats & score & beats \\
\midrule
displacement (CPA-class) & $1.828\pm0.213$ & --- & $2.309\pm0.223$ & --- \\
field composition        & $0.833\pm0.068$ & $3/3$ & $2.010\pm0.178$ & $3/3$ \\
$+$ state gate           & $0.633\pm0.037$ & $3/3$ & $\mathbf{0.700\pm0.079}$ & $3/3$ \\
$+$ covariant coupling   & $\mathbf{0.589\pm0.036}$ & $3/3$ & $0.772\pm0.081$ & $3/3$ \\
\bottomrule
\end{tabular}
\end{table}

\paragraph{A numerical failure worth reporting.} In a first three-seed run one organoid
seed with the coupling active returned NaN. The cause was not the model but the training
loop: gradient clipping does not stop a non-finite value, so a single non-finite forward
pass followed by an optimiser step wrote NaN into every parameter. We added two guards --
skip any step whose loss or gradient norm is non-finite (and report the count), and a
radial trust region on the composed field, active only on pathological early iterates. The
guarded rerun recovers that seed at $0.821$ and leaves every healthy run bit-identical
(finance seeds reproduce to four decimals). We note the trust region is a norm-based clip
and therefore \emph{not} itself covariant; it is inactive for trained models, where field
norms are $\sim\!0.08$ against a limit $50\times$ the state norm, and the covariance and
semigroup properties are stated for that regime.

\paragraph{Reading scores above 1.0.} On the organoid data an oracle that fits one mean
shift per condition scores $1.096$ on the \emph{training} split -- worse than inaction --
because populations are different patients and the same treatment moves them in directions
with mean pairwise cosine only $+0.23$ to $+0.35$. A model at $0.98$ on that split is
therefore beating the best possible fixed-shift predictor. The held-out splits are where
the comparison is meaningful.

\paragraph{Negative results.} (i) A \emph{global scalar} contraction does not help in
either domain ($0.882$ vs $0.862$ in finance, $1.965$ vs $1.947$ on the triple); state
dependence is what matters. (ii) On the finance exposure axis every model scores $6$--$9$,
and this is a property of the data, not the architecture: an oracle mean shift fitted on
the high-exposure split scores $0.287$ but the same oracle transferred from the training
split scores $2.788$, so that axis is a distribution shift rather than an extrapolation.
(iii) Saturation helps only where there is exposure range to extrapolate over; it is
neutral-to-harmful on the organoid ladder, whose four dose levels span no extrapolation.
(iv) Conditioning the operator on an encoding of the pre-intervention population -- the
mechanism of Meta Flow Matching -- improves the training split slightly ($1.092\to1.080$)
but damages the held-out triple ($0.903\to1.283$); with $26$ held-out populations the
encoder memorises replicate identity. We keep it in the codebase, off by default.

\section{Related work and limitations}

Flow matching on the Wasserstein manifold conditioned on a population is the contribution
of Meta Flow Matching \citep{atanackovic2024mfm}; we use that conditioning idea and do not
claim it. What is new here is the \emph{composition} of several generators with contraction
and coupling terms, which lets unseen intervention sets and exposures be predicted; MFM
conditions one velocity field and does not compose interventions. Metric-aware flow
matching \citep{kapusniak2024mfm} and energy-guided geometric flow matching
\citep{zweig2025eggfm} curve the interpolating \emph{paths} using a learned metric; we use a
metric for a different purpose -- to make a composition term tensorial -- and their metrics
are conformal rather than anisotropic. Learned intervention fields with Lie-bracket
geometry appear in recent causal-discovery work, restricted by its own statement of scope to
single-node interventions.

\paragraph{Limitations.} The covariance defect is measured on analytic diffeomorphisms; a
general proof is not given. The coupling term shows no gain over the contraction gate on
the organoid triple, where $26$ populations and $k{=}3$ leave it unidentifiable, so its
support rests on the finance composition split and on the covariance argument. Both domains
are evaluated in low-dimensional latent spaces ($d{=}6$, $d{=}8$), chosen because the
identifiability of a state-dependent second-order term degrades quickly with dimension.

\bibliographystyle{plainnat}
\begin{thebibliography}{9}
\bibitem[Atanackovic et al.(2025)]{atanackovic2024mfm}
L.~Atanackovic, X.~Zhang, B.~Amos, M.~Blanchette, L.~J.~Lee, Y.~Bengio, A.~Tong, and
K.~Neklyudov.
\newblock Meta Flow Matching: Integrating vector fields on the Wasserstein manifold.
\newblock \emph{ICLR}, 2025. arXiv:2408.14608.

\bibitem[Kapusniak et al.(2024)]{kapusniak2024mfm}
K.~Kapusniak et al.
\newblock Metric Flow Matching for smooth interpolations on the data manifold.
\newblock \emph{NeurIPS}, 2024. arXiv:2405.14780.

\bibitem[Ramos Zapatero et al.(2023)]{ramoszapatero2023}
M.~Ramos Zapatero et al.
\newblock Trellis tree-based analysis reveals stromal regulation of patient-derived
organoid drug responses.
\newblock \emph{Cell}, 186(25), 2023. doi:10.1016/j.cell.2023.11.005.

\bibitem[Zweig et al.(2025)]{zweig2025eggfm}
A.~Zweig, M.~Zhang, E.~Azizi, and D.~Knowles.
\newblock Energy Guided Geometric Flow Matching.
\newblock arXiv:2509.25230, 2025.
\end{thebibliography}

\end{document}
